\documentclass[12pt]{article}
\usepackage[utf8]{inputenc}
\usepackage[T1]{fontenc}
\usepackage[margin=1in]{geometry}
\usepackage{setspace}
\usepackage{xcolor}
\usepackage[colorlinks=true, urlcolor=blue!60!black, linkcolor=blue!60!black]{hyperref}
\setlength{\parindent}{0em}
\setlength{\parskip}{0.45em}
\renewcommand{\baselinestretch}{1.0}
\begin{document}
\begin{center}
{\Large \textit{Chaplinsky v. New Hampshire}, 315 U.S. 568 (1942)\par}
\vspace{0.3em}
{\normalsize Case Brief --- One-Page Edition\par}
\end{center}

\textbf{Memory Jogger.} Fighting words are unprotected: ``damned racketeer'' and ``damned Fascist,'' said to a marshal's face (\href{https://supreme.justia.com/cases/federal/us/315/568/}{Justia}).

\textbf{Facts.} Walter Chaplinsky, a Jehovah's Witness, denounced religion as a ``racket'' on a Rochester, New Hampshire, street. A crowd grew restless, and while being escorted toward the police station he called City Marshal Bowering ``a God damned racketeer'' and ``a damned Fascist.''

\textbf{Procedural History.} Convicted under a law barring ``offensive, derisive or annoying'' words in public, he lost in the New Hampshire Supreme Court, which limited the law to face-to-face fighting words. The Supreme Court affirmed.

\textbf{Votes.} 9--0: Murphy for a unanimous Court.

\textbf{Holding.} Narrow classes of speech, including fighting words, are unprotected, and the narrowly construed statute is neither vague nor an abridgment of free speech.

\textbf{Reasoning.} Free speech is ``not absolute at all times and under all circumstances.'' Obscenity, profanity, libel, and fighting words are ``no essential part of any exposition of ideas.'' The objective test asks what ``men of common intelligence'' would see as likely to provoke a fight. Later cases (\textit{Cohen}, \textit{Gooding}, \textit{R.A.V.}, \textit{Stevens}) sharply narrowed the doctrine.

\textbf{Examples.} \emph{Same-Side:} slurs and a dare to fight at a neighbor's door. \emph{Opposite-Side:} insults shouted at police by bullhorn. \emph{Fence-Sitter:} Chaplinsky's words spoken to an officer today.

\textbf{Critique.} His words were political criticism of an official, now protected.

\textbf{Key Quotation.} ``These include the lewd and obscene, the profane, the libelous, and the insulting or `fighting' words---those which by their very utterance inflict injury or tend to incite an immediate breach of the peace. It has been well observed that such utterances are no essential part of any exposition of ideas, and are of such slight social value as a step to truth that any benefit that may be derived from them is clearly outweighed by the social interest in order and morality.'' --- Justice Murphy, 315 U.S. at 572.

\end{document}
